\documentclass{ctexart}
\usepackage{Note}
\setCJKmainfont[
  BoldFont = 思源宋体 Bold,
  ItalicFont = 楷体,
]{宋体}
\author{2400011785 蒋锦豪}
\title{\tbf{结构化学\ 第3章作业}}
\begin{document}
\maketitle
\begin{homework}[1.]
    试证明BO势能面对于原子核构型的平动和转动具有不变性.
\end{homework}
\begin{homework}[2.]
    取键长$R=\SI{2}{\bohr}$,利用书中的公式分别画出$\ce{H_2^+}$的$1\sigma$成键/反键轨道对应的电子密度沿键轴方向的图像,并与相同几何构型下两个不成键的\ce{H}原子的1s轨道密度之和的一半进行对比,据此说明成键/反键轨道的电子密度相较于独立原子发生哪些变化.
\end{homework}
\begin{homework}[3.]
    取键轴为$z$方向,试分别估计\ce{F2}(键长约$\SI{135}{\pico\meter}$),乙炔(\ce{C#C}键长约$\SI{106}{\pico\meter}$)分子中两个F/C原子(1)$2\text{p}_z$轨道的重叠积分值; (2)$2\text{p}_x$轨道的重叠积分值.通过计算结果,你能得出什么结论?
\end{homework}
\begin{homework}[4.]
    在休克尔近似下,试编写程序计算并$N$苯($N=1,2,\cdots,6$,即苯,萘, $\cdots$, 并六苯)的HOMO和LUMO轨道能量(用$\alpha$和$\beta$表示),并作图表示趋势.
\end{homework}
\begin{homework}[5.]
    假定一个分子轨道具有$0.145A+0.844B$(非归一化)形式,找出与该分子轨道正交的轨道$A$和$B$的线性组合,并计算$S=0.250$时两种线性组合的归一化常数.
\end{homework}
\begin{homework}[6.]
    回答下面的问题.
    \begin{enumerate}[label=\tbf{(\alph*)}]
        \item 指出由p原子轨道肩并肩重叠形成的$\pi$成键与反键分子轨道的对称性特征(g或u).
        \item 指出由p原子轨道面对面重叠形成的$\delta$成键与反键分子轨道的对称性特征(g或u).
    \end{enumerate}
\end{homework}
\begin{homework}[8.]
    量子力学中的多体问题有如下形式的哈密顿量:
    \[\hat{H}=\sum_{k=1}^{N}\left[-\dfrac{1}{2}\dfrac{\p^2}{\p x_k^2}+v(x_k)\right]+\sum_{1\leq k<l\leq N}w(x_k,x_l)\]
    现在我们考虑一个势能项均为简谐势的一维体系:
    \[v(x_k)=\dfrac{1}{2}x_k^2,\quad w(x_k,x_l)=\dfrac{\lambda}{2}(x_k-x_l)^2\]
    \begin{enumerate}[label=\tbf{(\alph*)}]
        \item 假定粒子没有自旋,证明此体系具有以下形式的基态波函数与基态能量:
        \[\iPsi_0=C_0\exp\left\{-\dfrac12\sum_{k=1}^{N}x_k^2-\dfrac{\theta}{2}\sum_{1\leq k<l\leq N}(x_k-x_l)^2\right\}\]
        \[E_0=\dfrac{1+(N-1)(N\theta+1)}{2}\]
        其中
        \[C_0=\dfrac{(N\theta+1)^{(N-1)/4}}{\pi^{N/4}},\quad\theta=\dfrac{\sqrt{\lambda N+1}-1}{N}\]
        提示:可以寻找一组新的坐标$\{q_k\}_{k=1}^{N}$使得$\hat{H}$具有如下形式:
        \[\hat{H}=\sum_{k=1}^{N}\left[-\dfrac{1}{2}\dfrac{\p^2}{\p q_k^2}+\dfrac{1}{2}\omega_k^2 q_k^2\right]\]
        \item 观察可知, \tbf{(a)}中的$\iPsi_0$对应于玻色子的基态.试证明以下波函数给出体系的一个费米子本征态:
        \[\iPsi_F=C_F\begin{vmatrix}
            1&1&\cdots&1\\
            x_1&x_2&\cdots&x_N\\
            \vdots&\vdots&\ddots&\vdots\\
            x_1^{N-1}&x_2^{N-1}&\cdots&x_{N}^{N-1}
        \end{vmatrix}\iPsi_0\]
        其中$C_F$为归一化系数.实际上$\iPsi_F$就是体系的费米子基态(以下称为费米子的基态正解),其能量为
        \[E_F=\dfrac{1+(N^2-1)(N\theta+1)}{2}\]
        \item 请写出此体系的Hartree-Fock方程(不考虑自旋).
        \item 假设我们考虑的是$N=m+n$个有自旋的电子,其中$m$个是$\alpha$电子, $n$个是$\beta$电子,证明体系有如下本征态:
        \[\iPsi=C\begin{vmatrix}
            \alpha(1)&\alpha(2)&\cdots&\alpha(N)\\
            x_1\alpha(1)&x_2\alpha(2)&\cdots&x_N\alpha(N)\\
            \vdots&\vdots&\ddots&\vdots\\
            x_1^{m-1}\alpha(1)&x_2^{m-1}\alpha(2)&\cdots&x_N^{m-1}\alpha(N)\\
            \beta(1)&\beta(2)&\cdots&\beta(N)\\
            x_1\beta(1)&x_2\beta(2)&\cdots&x_N\beta(N)\\
            \vdots&\vdots&\ddots&\vdots\\
            x_1^{n-1}\beta(1)&x_2^{n-1}\beta(2)&\cdots&x_N^{n-1}\beta(N)
        \end{vmatrix}\iPsi_0\]
        这里$C$是归一化系数, $\alpha(k)$表示第$k$个标号的电子是$\alpha$电子, $\beta(l)$表示第$l$个标号的电子为$\beta$电子. $\iPsi$对应的本征能量为
        \[E=\dfrac{1+(m^2+n^2-1)(N\theta+1)}{2}\]
        \item 证明在\tbf{(d)}中,如果$N$为偶数,则体系的能量在$m=n$时达到最低(单重态);如果$N$为奇数,则体系的能量在$|m-n|=1$时达到最低(二重态).
        \item 证明对于\tbf{(d)}中$N=3$的最稳定的状态,总波函数不能写成空间波函数与自旋波函数的乘积:
        \[\iPsi(x_1,\sigma_1,x_2,\sigma_2,x_3,\sigma_3)\neq\Phi(x_1,x_2,x_3)\chi(\sigma_1,\sigma_2,\sigma_3)\]
        这里$\sigma_i\in{\alpha,\beta}$是第$i$个电子的自旋维度.
        \item 在考虑自旋时,证明$N=2$时Hartree-Fock方程给出如下基态解:
        \[\iPsi^{\text{HF}}=\dfrac{1}{\sqrt2}\psi(x_1)\psi(x_2)\left[\alpha(1)\beta(2)-\alpha(2)\beta(1)\right]\]
        其中
        \[\psi(x)=\left(\dfrac{1+\lambda}{\pi^2}\right)^{1/8}\exp\left(-\dfrac{\sqrt{1+\lambda}}{2}x^2\right)\]
        \item 考虑\tbf{(g)}中的情形,通过计算分别写出正解和Hartree-Fock解的电子密度表达式,并作图比较$\lambda=-\dfrac49,-\dfrac29,0,1,5,10$的时的情形.当$\lambda\to+\infty$时,两者有何区别?当$\lambda\to-\dfrac12$时,两者有何区别?
    \end{enumerate}
\end{homework}
\begin{solution}
\begin{enumerate}[label=\tbf{(\alph*)}]
    \item 考虑线性空间$V$中的向量$\vec{x}=(x_1,x_2,\cdots,x_N)$,有
    \[\begin{aligned}
        \vec{x}^\t\mat{V}\vec{x}
        &=\sum_{k=1}^{N}v(x_k)+\sum_{1\leq k<l\leq N}w(x_k,x_l)\\
        &=\dfrac12\sum_{k=1}^{N}x_k^2+\dfrac{\lambda}{2}\sum_{1\leq k<l\leq N}\left(x_k-x_l\right)^2\\
        &=\dfrac12\sum_{k=1}^{N}x_k^2+\dfrac{\lambda}{2}\left[(N-1)\sum_{k=1}^{N}x_k^2-2\sum_{1\leq k<l\leq N}x_kx_l\right]\\
        &=\dfrac{\lambda(N-1)+1}{2}\sum_{k=1}^{N}x_k^2-\lambda\sum_{1\leq k<l\leq N}x_kx_l
    \end{aligned}\]
    其中
    \[\mat{V}=\dfrac12(\lambda N+1)\mat{I}-\dfrac\lambda2\mat{E}\]
    现在考虑$\mat{V}$的所有特征值.首先注意到对任意向量$\vec{u}$都有$\mat{E}\vec{u}=\displaystyle\left(\sum_{k=1}^{N}u_k\right)(1,1,\cdots,1)$.因此令$\vec{q}_1=\dfrac{1}{\sqrt{N}}(1,1,\cdots,1)$则有
    \[\mat{V}\vec{q}_1=\dfrac{\lambda N+1}{2}\vec{q}_1-\dfrac{\lambda N}{2}\vec{q}_1=\dfrac{1}{2}\vec{q}_1\]
    除此之外,所有满足$\displaystyle\sum_{k=1}^{N}u_k=0$的非零向量$\vec{u}$都是$\vec{V}$的特征向量,对应的特征值为$\dfrac{\lambda N+1}{2}$.所有这样的$\vec{u}$构成的线性空间$U$的维数为$N-1$,因此该特征值的重数为$N-1$.取$U$的一组规范正交基$\{\vec{q}_k\}_{k=2}^{N}$,这样$\{\vec{q}_k\}_{k=1}^{N}$是$V$的一组规范正交基.于是令
    \[\mat{P}=\begin{bmatrix}
        \vec{q}_1&\vec{q}_2&\cdots&\vec{q}_N
    \end{bmatrix}\]
    则有
    \[\mat{P}^\t\mat{V}\mat{P}=\diag\left\{\dfrac{1}{2},\dfrac{1+\lambda N}{2},\cdots,\dfrac{1+\lambda N}{2}\right\}\]
    于是可以令
    \[\vec{q}=\mat{P}\vec{x},\quad\vec{x}^\t\mat{V}\vec{x}=\vec{q}^\t\diag\left\{\dfrac{1}{2},\dfrac{1+\lambda N}{2},\cdots,\dfrac{1+\lambda N}{2}\right\}\vec{q}=\dfrac12q_1^2+\dfrac{1+\lambda N}{2}\sum_{k=2}^{N}q_k^2\]
    令$\omega_1=1$, $\omega_k=\sqrt{\lambda N+1}$(其中$k=2,\cdots,N$)则有
    \[\sum_{k=1}^{N}v(x_k)+\sum_{1\leq k<l\leq N}w(x_k,x_l)=\sum_{k=1}^{N}\omega_k^2q_k^2\]
    由于$\vec{q}=\mat{P}\vec{x}$是规范正交变换,因此Laplace算子的形式相同,即
    \[\sum_{k=1}^{N}\dfrac{\p^2}{\p x_k^2}=\sum_{k=1}^{N}\dfrac{\p^2}{\p q_k^2}\]
    于是体系的薛定谔方程可以改写为
    \[\hat{H}\iPsi(q_1,\cdots,q_N)=E\iPsi(q_1,\cdots,q_N),\quad\hat{H}=\sum_{k=1}^{N}\left[-\dfrac{1}{2}\dfrac{\p^2}{\p q_k^2}+\dfrac{1}{2}\omega_k^2 q_k^2\right]\]
    分离变量
    \[\iPsi(q_1,\cdots,q_N)=\prod_{k=1}^{N}\psi_k(q_k)\]
    代入薛定谔方程并整理可得
    \[-\dfrac{1}{2}\dfrac{\p\psi_k}{\p q_k^2}+\dfrac12\omega_k^2q_k^2\psi_k=E_k\psi_k,\quad k=1,2,\cdots,N\]
    这是典型的谐振子模型,其基态解为
    \[\psi_{k,0}(q_k)=\left(\dfrac{\omega_k}{\pi}\right)^{1/4}\exp\left(-\dfrac{\omega_k}{2}q_k^2\right),\quad E_k=\dfrac{1}{2}\omega_k\]
    于是
    \[\iPsi_0=\prod_{k=1}^{N}\psi_k(q_k)=\dfrac{\omega_1^{1/4}\prod_{k=2}^{N}\omega_k^{1/4}}{\pi^{N/4}}\exp\left\{-\dfrac12\sum_{k=1}^{N}\omega_kq_k^2\right\}\]
    而
    \[\omega_1=1,\quad\omega_k=\sqrt{\lambda N+1}=N\theta+1,\quad k=2,\cdots,N\]
    于是
    \[C_0=\dfrac{\omega_1^{1/4}\prod_{k=2}^{N}\omega_k^{1/4}}{\pi^{N/4}}=\dfrac{(N\theta+1)^{(N-1)/4}}{\pi^{N/4}}\]
    根据前面构造的$\vec{q}_1$可得$\displaystyle q_1=\dfrac{x_1+\cdots+x_N}{\sqrt{N}}$,而
    \[\vec{q}^\t\vec{q}=\vec{x}^\t\mat{P}^\t\mat{P}\vec{x}=\vec{x}^\t\mat{I}\vec{x}=\vec{x}^\t\vec{x}\]
    于是
    \[\sum_{k=2}^{N}q_k^2=\sum_{k=1}^{N}x_k^2-q_1^2=\sum_{k=1}^{N}x_k^2-\dfrac{1}{N}\left(\sum_{k=1}^{N}x_k\right)^2=\dfrac{1}{N}\sum_{1\leq k<l\leq N}(x_k-x_l)^2\]
    于是
    \[\sum_{k=1}^{N}\omega_kq_k^2=\dfrac{1}{N}\left(\sum_{k=1}^{N}x_k\right)^2+\dfrac{\sqrt{\lambda N+1}}{N}\sum_{1\leq k<l\leq N}(x_k-x_l)^2=\sum_{k=1}^{N}x_k^2+\dfrac{\sqrt{\lambda N+1}-1}{N}\sum_{1\leq k<l\leq N}(x_k-x_l)^2\]
    于是
    \[\iPsi_0=C_0\exp\left\{-\dfrac12\sum_{k=1}^{N}\omega_kq_k^2\right\}=C_0\exp\left\{-\dfrac12\sum_{k=1}^{N}x_k^2-\dfrac{\theta}{2}\sum_{1\leq k<l\leq N}(x_k-x_l)^2\right\}\]
    基态能量则为
    \[E_0=\dfrac12\sum_{k=1}^{N}\omega_k=\dfrac{1+(N-1)(N\theta+1)}{2}\]
    证毕.
    \item 在\tbf{(a)}的基础上,费米子构成的波函数应当满足交换反对称性,即交换$x_k$与$x_l$后波函数反号.根据行列式的性质,交换其中的两列则行列式反号,而根据$\iPsi_0$的定义可知$\iPsi_0$不变,于是即$\iPsi_F$变为$\iPsi_F$,满足交换反对称性.\\
    现在证明$\iPsi_F$仍为$\hat{H}$的本征态.根据范德蒙德行列式的性质,令
    \[D(x_1,\cdots,x_N)=\begin{vmatrix}
        1&1&\cdots&1\\
        x_1&x_2&\cdots&x_N\\
        \vdots&\vdots&\ddots&\vdots\\
        x_1^{N-1}&x_2^{N-1}&\cdots&x_{N}^{N-1}
    \end{vmatrix}=\prod_{i<j}(x_i-x_j)\]
    于是
    \[\begin{aligned}
        \hat{H}(D\iPsi_0)
        &=\left(-\dfrac{1}{2}\sum_{k=1}^{N}\p_k^2+V\right)(D\iPsi_0)\\
        &=-\dfrac12\sum_{k=1}^{N}\left(\iPsi_0\p_k^2 D+2\p_k\iPsi_0\p_k D+D\p_k^2\iPsi_0\right)+VD\iPsi_0\\
        &=-\dfrac12\iPsi_0\sum_{k=1}^{N}\left(\p_k^2 D+2\p_k\ln\iPsi_0\p_k D\right)+D\underbrace{\left(-\dfrac{1}{2}\sum_{k=1}^{N}\p_k^2+V\right)\iPsi_0}_{\hat{H}\iPsi_0}\\
        &=\iPsi_0\left[-\dfrac12\sum_{k=1}^{N}\left(\p_k^2+2\p_k\ln\iPsi_0\p_k \right)\right]D+E(D\iPsi_0)
    \end{aligned}\]
    记
    \[\hat{\mathcal{L}}=-\dfrac12\sum_{k=1}^{N}\left(\p_k^2+2\p_k\ln\iPsi_0\p_k\right)\]
    只需证明$D$是$\hat{\mathcal{L}}$的本征态.首先
    \[\p_k\ln\iPsi_0=\p_k\left\{-\dfrac12\sum_{k=1}^{N}x_k^2-\dfrac{\theta}{2}\sum_{1\leq k<l\leq N}(x_k-x_l)^2\right\}=-x_k-\theta\sum_{k\neq l}(x_k-x_l)\]
    注意到
    \[\ln D=\sum_{k<l}\ln(x_l-x_k)\]
    于是
    \[\p_k D=D\p_k\ln D=D\sum_{k\neq l}\dfrac{1}{x_k-x_l}\]
    \[\p_k^2 D=D\left[\left(\sum_{k\neq l}\dfrac{1}{x_k-x_l}\right)^2-\sum_{k\neq l}\dfrac{1}{(x_k-x_l)^2}\right]\]
    而
    \[\sum_{k=1}^{N}\p_k^2 D=D\sum_{k=1}^{N}\sum_{l\neq k}\sum_{j\neq k}\dfrac{2}{(x_k-x_l)(x_k-x_j)}=2D\sum_{j\neq k\neq l}\dfrac{(x_j-x_l)+(x_l-x_k)+(x_k-x_j)}{(x_k-x_l)(x_l-x_j)(x_j-x_k)}=0\]
    于是二阶偏导项为零,只需考虑一阶偏导项.如此则有
    \[\begin{aligned}
        \hat{\mathcal{L}}D
        &=-\dfrac12\sum_{k=1}^{N}\left\{2\left[-x_k-\theta\sum_{k\neq l}(x_k-x_l)\right]\sum_{k\neq l}\dfrac{1}{x_k-x_l}\right\}\\
        &=D\left[(1+N\theta)\sum_{k=1}^{N}\sum_{l\neq k}\dfrac{x_k}{x_k-x_l}+\sum_{j=1}^{N}x_j\sum_{k=1}^{N}\sum_{l\neq k}\dfrac{1}{x_k-x_l}\right]\\
        &=D\left[\dfrac{(1+N\theta)N(N-1)}{2}+0\right]\\
        &=\dfrac{(1+N\theta)N(N-1)}{2}D
    \end{aligned}\]
    于是$D$确实是$\hat{\mathcal{L}}$的本征态.这样就有
    \[\hat{H}(D\iPsi_0)=\left[\dfrac{(1+N\theta)N(N-1)}{2}+E_0\right](D\iPsi_0)\]
    于是$\iPsi_F$是$\hat{H}$的本征态,其本征值为
    \[E_F=\dfrac{(1+N\theta)N(N+1)}{2}+E_0\dfrac{1+(N-1)(N\theta+1)+(1+N\theta)N(N-1)}{2}=\dfrac{1+(N^2-1)(N\theta+1)}{2}\]
    证毕.
    \item 体系的Hartree-Fock方程为
    \item 注意到题中行列式的第$k$列可以写作
    \[\begin{bmatrix}
        \alpha(k)\\\vdots\\x_{k}^{m-1}\alpha(k)\\\beta(k)\\\vdots\\x_{k}^{n-1}\beta(k)
    \end{bmatrix}=\alpha(k)\begin{bmatrix}
        1\\\vdots\\x_{k}^{m-1}\\0\\\vdots\\0
    \end{bmatrix}+\beta(k)\begin{bmatrix}
        0\\\vdots\\0\\1\\\vdots\\x_k^{n-1}
    \end{bmatrix}=\alpha(k)\vec{a}_k+\beta(k)\vec{b}_k\]
    考虑行列式的每一列按上述方法拆分后的求和项.由于$\vec{a}_k$只在前$m$行非零, $\vec{b}_k$只在后$n$行非零;而根据行列式的定义,必须在前$m$行选择$m$个不同列的元素,后$n$行选择$n$个不同列的元素构成求和项,因此必须选择$m$列$\vec{a}$向量和$n$列$\vec{b}$向量才能使得求和项不全为$0$.于是
    \[|\cdot|=\sum_{k_1k_2\cdots k_N}\prod_{i=1}^{m}\alpha\left(k_i\right)\prod_{j=m+1}^{N}\beta\left(k_j\right)\det\mat{M}_{k_1\cdots k_N}\]
    其中$k_1k_2\cdots k_N$是$1$到$N$的一个排列,满足$k_1<\cdots<k_m$, $k_{m+1}<\cdots<k_N$. $\mat{M}$是由$\vec{a}_{k_1},\cdots,\vec{a}_{k_m}$和$\vec{b_{k_{m+1}}},\cdots,\vec{b}_{k_N}$按照列号大小组成的矩阵.而
    \[\begin{aligned}
        \left|\det\mat{M}\right|
        &=\begin{vmatrix}
        \vec{a}_{k_1}&\cdots&\vec{a}_{k_m}&\vec{b_{k_{m+1}}}&\cdots&\vec{b}_{k_N}
    \end{vmatrix}\\
    &=\begin{vmatrix}
        1&\cdots&1\\
        \vdots&\ddots&\vdots\\
        x_{k_1}^{m}&\cdots&x_{k_m}^{m}
    \end{vmatrix}\begin{vmatrix}
        1&\cdots&1\\
        \vdots&\ddots&\vdots\\
        x_{k_{m+1}}^{n}&\cdots&x_{k_N}^{n}
    \end{vmatrix}\\
    &=\prod_{1\leq i<j\leq m}\left(x_{k_j}-x_{k_i}\right)\prod_{m<i<j\leq N}\left(x_{k_j}-x_{k_i}\right)
    \end{aligned}\]
    现在对于符合前设的所有排列$k_1\cdots k_N$, 考虑任意一项
    \[D_{k_1\cdots k_N}=(\pm)\prod_{i=1}^{m}\alpha\left(k_i\right)\prod_{j=m+1}^{N}\beta\left(k_j\right)\underbrace{\prod_{1\leq i<j\leq m}\left(x_{k_j}-x_{k_i}\right)}_{D_\alpha}\underbrace{\prod_{m<i<j\leq N}\left(x_{k_j}-x_{k_i}\right)}_{D_\beta}\]
    则有
    \[\hat{H}\left(D_{k_1\cdots k_N}\iPsi_0\right)=\iPsi_0\hat{\mathcal{L}}D_{k_1\cdots k_N}+E_0\left(D_{k_1\cdots k_N}\iPsi_0\right)\]
    观察到$D_{k_1\cdots k_N}$由两部分不相关的项构成.令
    \[\p_\alpha^2=\sum_{i=1}^{m}\p^2_{x_{k_i}},\quad\p_\beta^2=\sum_{j=m+1}^{N}\p^2_{x_{k_j}},\quad D=D_\alpha D_\beta\]
    于是
    \[\begin{aligned}
        \hat{\mathcal{L}}D
        &=-\dfrac12\sum_{k=1}^{N}\left(\p_k^2+2\p_k\ln\iPsi_0\p_k\right)D\\
        &=-\dfrac12\sum_{i=1}^{m}\left(\p_{k_i}^2+2\p_{k_i}\ln\iPsi_0\p_{k_i}\right)D-\dfrac12\sum_{j=m+1}^{N}\left(\p_{k_j}^2+2\p_{k_j}\ln\iPsi_0\p_{k_j}\right)D\\
        &=-\dfrac12\left[D_\beta\left(\p_\alpha^2+2\sum_{i=1}^{m}\p_{k_i}\ln\iPsi_0\p_{k_i}\right)D_\alpha+D_\alpha\left(\p_\beta^2+2\sum_{j=m+1}^{N}\p_{k_j}\ln\iPsi_0\p_{k_j}\right)D_\beta\right]
    \end{aligned}\]
    类似前一问可得
    \[-\left(\p_\alpha^2+\sum_{i=1}^{m}\p_{k_i}\ln\iPsi_0\p_{k_i}\right)D_\alpha=\dfrac{(1+N\theta)m(m-1)}{2}D_\alpha\]
    \[-\left(\p_\beta^2+\sum_{j=m+1}^{N}\p_{k_j}\ln\iPsi_0\p_{k_j}\right)D_\beta=\dfrac{(1+N\theta)n(n-1)}{2}D_\beta\]
    这样就有
    \[\hat{L}D=\dfrac{(1+N\theta)(m^2+n^2-m-n)}{2}D_\alpha D_\beta=\dfrac{(1+N\theta)(m^2+n^2-m-n)}{2}D\]
    因此原行列式的每一项与$\iPsi_0$的乘积都是$\hat{H}$的本征态.于是
    \[\begin{aligned}
        \hat{H}(|\cdot|\iPsi_0)
        &=\sum_{k_1\cdots k_N}\hat{H}\left[(\pm)\prod_{i=1}^{m}\alpha\left(k_i\right)\prod_{j=m+1}^{N}\beta\left(k_j\right)D_\alpha D_\beta\iPsi_0\right]\\
        &=\sum_{k_1\cdots k_N}\left[E_0+\dfrac{(1+N\theta)(m^2+n^2-N)}{2}\right]\left[(\pm)\prod_{i=1}^{m}\alpha\left(k_i\right)\prod_{j=m+1}^{N}\beta\left(k_j\right)D_\alpha D_\beta\iPsi_0\right]\\
        &=\dfrac{1+(m^2+n^2-1)(N\theta+1)}{2}\sum_{k_1\cdots k_N}\left[(\pm)\prod_{i=1}^{m}\alpha\left(k_i\right)\prod_{j=m+1}^{N}\beta\left(k_j\right)D_\alpha D_\beta\iPsi_0\right]\\
        &=\dfrac{1+(m^2+n^2-1)(N\theta+1)}{2}\sum_{k_1\cdots k_N}|\cdot|\iPsi_0
    \end{aligned}\]
    于是$\iPsi=|\cdot|\iPsi_0$确实是$\hat{H}$的本征态,对应的本征能量为
    \[E=\dfrac{1+(m^2+n^2-1)(N\theta+1)}{2}\]
    证毕.
    \item 将$n=N-m$代入$E$中有
    \[E=\dfrac12+\dfrac{N\theta+1}{2}\left[m^2+(N-m)^2\right]=\dfrac12+\dfrac{N\theta+1}{2}(2m^2-2mN+N^2)\]
    由于$N\theta+1>0$,因此当$E$最小时$2m^2-2mN+N^2$最小.这二次函数开口向上,对称轴为$m_0=\dfrac N2$.当$N$为偶数时,取$m=n=\dfrac N2$使$E$最小;当$N$为奇数时,取$m=\dfrac{N\pm1}{2}$,也即$|m-n|=1$时使$E$最小.
    \item $N=3$的最稳定状态为$m=2$, $n=1$或$m=1$, $n=2$.不妨考虑前者,此时
    \[\iPsi=C\begin{vmatrix}
        \alpha(1)&\alpha(2)&\alpha(3)\\
        x_1\alpha(1)&x_2\alpha(2)&x_3\alpha(3)\\
        \beta(1)&\beta(2)&\beta(3)
    \end{vmatrix}\iPsi_0\]
    \item 1
    \item 正解为
    \[\iPsi=C\begin{vmatrix}
        \alpha(1)&\alpha(2)\\
        \beta(1)&\beta(2)
    \end{vmatrix}\iPsi_0=C\begin{vmatrix}
        \alpha(1)&\alpha(2)\\
        \beta(1)&\beta(2)
    \end{vmatrix}C_0\exp\left[-\dfrac12(x_1^2+x_2^2)-\dfrac{\theta}{2}(x_1-x_2)^2\right]\]
    整理并归一化可得
    \[\iPsi=\dfrac{(2\lambda+1)^{1/8}}{(2\pi)^{1/2}}\exp\left[-\dfrac12(x_1^2+x_2^2)-\dfrac{\sqrt{2\lambda+1}-1}{4}(x_1-x_2)^2\right]\left[\alpha(1)\beta(2)-\alpha(2)\beta(1)\right]\]
    于是两者的空间波函数分别为
    \[\psi^{\text{HF}}(x_1,x_2)=\left(\dfrac{1+\lambda}{\pi^2}\right)^{1/4}\exp\left[-\dfrac{\sqrt{1+\lambda}}{2}\left(x_1^2+x_2^2\right)\right]\]
    \[\psi(x_1,x_2)=\left(\dfrac{2\lambda+1}{\pi^4}\right)^{1/8}\exp\left[-\dfrac12(x_1^2+x_2^2)-\dfrac{\sqrt{2\lambda+1}-1}{4}(x_1-x_2)^2\right]\]
    电子密度为空间波函数的平方,即
    \[\rho^{\text{HF}}(x_1,x_2)=\dfrac{\sqrt{1+\lambda}}{\pi}\exp\left[-\sqrt{1+\lambda}\left(x_1^2+x_2^2\right)\right]\]
    \[\rho(x_1,x_2)=\dfrac{(2\lambda+1)^{1/4}}{\pi}\exp\left[-\left(x_1^2+x_2^2\right)-\dfrac{\sqrt{2\lambda+1}-1}{2}(x_1-x_2)^2\right]\]
    $\lambda$取不同值时的电子密度示意图如下(蓝色为$\rho^{\text{HF}}$,红色为$\rho$):
    \begin{figure}[H]\centering
        \subfigure[$\lambda=-\dfrac49$]{
        \begin{minipage}[b]{.45\linewidth}
            \centering\includegraphics[scale=0.75]{figure/8h-1.eps}
        \end{minipage}}
        \subfigure[$\lambda=-\dfrac29$]{
        \begin{minipage}[b]{.45\linewidth}
            \centering\includegraphics[scale=0.75]{figure/8h-2.eps}
        \end{minipage}}
        \subfigure[$\lambda=0$(取值相同)]{
        \begin{minipage}[b]{.45\linewidth}
            \centering\includegraphics[scale=0.75]{figure/8h-3.eps}
        \end{minipage}}
        \subfigure[$\lambda=1$]{
        \begin{minipage}[b]{.45\linewidth}
            \centering\includegraphics[scale=0.75]{figure/8h-4.eps}
        \end{minipage}}
        \subfigure[$\lambda=5$]{
        \begin{minipage}[b]{.45\linewidth}
            \centering\includegraphics[scale=0.75]{figure/8h-5.eps}
        \end{minipage}}
        \subfigure[$\lambda=10$]{
        \begin{minipage}[b]{.45\linewidth}
            \centering\includegraphics[scale=0.75]{figure/8h-6.eps}
        \end{minipage}}
        \caption{$\lambda$取不同值时$\rho^{\text{HF}}$与$\rho$的比较}
    \end{figure}
    当$\lambda\to-\dfrac12$时,吸引势趋向极限,正解明显趋向于仅在$x_1=x_2$处$\rho\neq0$,而HF解仍给出独立的高斯分布;当$\lambda\to+\infty$时,排斥势趋向极限,正解明显趋向于在$x_1=x_2$处$\rho=0$,而HF解仍给出独立的高斯分布.因此, HF解不能描述粒子之间的相互作用.
\end{enumerate}
\end{solution}
\end{document}